\section{金属材料：結合・結晶・組織から試験と選択へ}
\label{theory:material:overview}
材料選択では、単に「強い金属」を探すだけでは足りない。小さな荷重での変形、永久変形の始まり、き裂の進展、繰返し負荷、電気抵抗、加工のしやすさは異なる特性である。組成と加工条件が結晶・転位・粒界・欠陥を変え、それが使用中の応答へ結び付く。本節は基礎科目の材料問題から、その因果関係を整理する。

\subsection{金属結合と塑性：電子と転位の役割}
\label{theory:material:bonding}
金属の外側の電子は原子一つに局在せず、正の原子核・内殻電子からなる芯の間で共有される。自由電子の海という近似は金属結合と導電性を理解する入口である\cite{mit_material_atomistic2001}。同位体は中性子数の違う同じ元素の原子であり、自由電子の代わりの結合機構ではない。

結合が特定方向だけに限られないことと、原子面が容易に全て同時に滑ることは同じではない。実際の塑性では格子の線欠陥である転位が移動し、少しずつすべりを進める。転位の移動を妨げる粒界・溶質・析出物は強化に寄与するが、延性との交換条件も生む\cite{mit_material_dislocations2001}。結合だけで全金属の室温延性を保証しない。

\subsection{BCC・FCC・HCP：胞と配位数を数え分ける}
\label{theory:material:crystals}
単位胞は結晶の並進で繰り返す領域である。格子点と、その各点に付く原子の基底を区別する。胞の形を決めずに「一胞に何原子」と尋ねると、原始胞と大きな慣用胞の数が混ざる\cite{cambridge_unitcell,princeton_hcp}。

通常の立方胞では、角の原子は隣り合う8胞で共有し、面の原子は2胞で共有する。この分率から独自に数えると、
\[
 N_{\rm BCC}=8\times\frac18+1=2,\qquad
 N_{\rm FCC}=8\times\frac18+6\times\frac12=4.
\]
BCCの最近接配位数は8、FCCは12である。配位数は一つの原子の周囲にある最近接原子の数であり、胞の中の原子数ではない\cite{openstax_crystals}。

硬球近似で半径を\(R\)、立方胞辺を\(a\)とする。BCCは体対角線、FCCは面対角線で接触するため、幾何から
\[
 \sqrt3a=4R\quad({\rm BCC}),\qquad
 \sqrt2a=4R\quad({\rm FCC})
\]
となる。FCCでは\(a=2\sqrt2R\)である。式の暗記だけでなく、接触する原子列を指定する。

HCPは六方Bravais格子に二原子基底を置いた構造であり、原始並進胞に2原子を含む。よく図示される大きな六角柱では、角12原子の共有分\(12/6\)、上下の面中心2原子の共有分\(2/2\)、内部3原子を合計し、6原子となる。この六角柱は原始並進胞の3倍の体積を持つ。平成29年度と令和6年度の「HCPは2原子」は前者の規約で読み、六角柱の図へ無条件に転記しない。

常温常圧の元素ではFe・NaはBCC、Al・CuはFCC、Ti・ZnはHCPである\cite{princeton_iron,princeton_sodium,princeton_aluminum,princeton_copper,princeton_titanium,princeton_zinc}。相の安定性は温度・圧力・組成に依存する。別資料の温度を指定しない元素例示を、室温の分類にそのまま使わない。

\subsection{荷重と変位から応力・ひずみへ}
\label{theory:material:tension}
初期断面積を\(A_0\)、初期標点距離を\(L_0\)、荷重を\(F\)、標点間の長さを\(L\)とする。公称値は
\[
 \sigma_{\rm e}=\frac{F}{A_0},\qquad
 \varepsilon_{\rm e}=\frac{L-L_0}{L_0}
\]
である。試験機のクロスヘッド変位は機械・つかみの変形を含むことがあるため、試験片の標点間の伸びと対応させる。

現在の断面積\(A\)で割った\(\sigma_{\rm t}=F/A\)が真応力である。真ひずみは現在長さによる微小増分を積分するので、
\[
 \varepsilon_{\rm t}=\int_{L_0}^{L}\frac{d\ell}{\ell}
 =\ln\frac{L}{L_0}=\ln(1+\varepsilon_{\rm e}).
\]
全伸びを変形後の長さで割るだけでは、この積分値にならない\cite{mit_material_stressstrain2001}。微小ひずみでは両ひずみは近い。大変形・くびれ後には、標点間平均と局所値を区別する必要がある。

\subsection{フック則：剛性と強度を分ける}
\label{theory:material:elasticity}
小応力の単軸線形弾性では\(\sigma=E\varepsilon\)となる。比例定数\(E\)がヤング率であり、ひずみが無次元なので応力と同じPaを単位に持つ。棒の剛性は\(k=EA_0/L_0\)であり、材料の\(E\)に加えて形状にも依存する\cite{mit_material_elastic2000}。

降伏応力は塑性が始まる応力、引張強度は通常、公称応力の最大値である。ヤング率はこの二つの強度ではない。ポアソン比は横と軸のひずみの比であり無次元である。ヘンリー則は溶解、ブラッグ則は回折の関係である。

出題のMg45 GPa、W407 GPaは大学教材にもある代表値である\cite{princeton_magnesium,princeton_tungsten}。同じ応力ならMgの方が大きく弾性変形する。通常の金属では昇温に伴う弾性率低下を考えるが、方位・温度区間・相変態を指定し、全材料の単調減少へ一般化しない。弾性は除荷による回復を指し、線形性と常に同義ではない。

\subsection{粒径・加工硬化・再結晶}
\label{theory:material:strengthening}
一般的な多結晶鋼などでは、結晶粒界が転位移動の障害になる。通常の粒径範囲の経験的関係は
\[
 \sigma_y=\sigma_0+\frac{k_y}{\sqrt d}
\]
である。\(d\)が小さいほど第二項が大きくなる。粗大化で強度が上がるという出題肢は、この通常の比較では逆である。

塑性加工で転位が増えると、転位同士も移動を妨げる。これが加工硬化を理解する基本である。加熱による回復・再結晶は蓄積した欠陥を減らし、通常は加工硬化を緩和して延性を回復する\cite{mit_material_dislocations2001}。

Hall--Petch式を原子数が極端に少ないナノ粒まで外挿し、強度が際限なく増すとはしない。粒界すべり・回転・拡散や加工時の欠陥など、支配機構が変わる。逆Hall--Petchの報告にも加工・硬度から降伏応力への換算・試験中粒成長の問題があり、全てのナノ材料が必ず弱化するという保証にしない\cite{hall_petch_limits2020}。

\subsection{欠陥・脆性破壊・疲労・衝撃靭性}
\label{theory:material:fracture}
無欠陥の原子結合を一斉に切る理論強度と、孔・介在物・き裂を含む実材料の強度は異なる。き裂先端には応力が集中し、き裂寸法・形状・負荷と破壊抵抗を一緒に評価する。線形破壊力学の典型的形は\(K_I=Y\sigma\sqrt{\pi a}\)であり、\(a\)の定義と形状係数\(Y\)を指定する。

脆性破壊は顕著な塑性変形が小さい破壊、疲労は繰返し負荷による損傷蓄積である。疲労では一回なら耐えられる荷重でも、き裂が成長して最後に急速破断することがある\cite{mit_material_fatigue2001}。腐食疲労はあるが、疲労を経年腐食だけで定義しない。

リバティ船とコメットは破壊力学発展を促した代表例である\cite{nist_fracture_history2010}。出題は船の脆性と航空機の疲労を対照する。コメットに関する大学教材は窓周辺のリベット孔からのき裂も説明する\cite{mit_material_fracture2001}。原典の「窓の角」という簡略表現を全事故の起点の保証へ変えず、溶接部・開口形状・低温・繰返し圧力などの複合条件を考える。

Charpyは切欠き試験片を衝撃破断し、吸収エネルギーから衝撃靭性を評価する。温度・寸法・切欠き・負荷速度に依存する測定であり、平面ひずみ破壊靭性\(K_{IC}\)の直接測定ではない\cite{nist_charpy}。

\subsection{純鉄の結晶相と磁気転移}
\label{theory:material:ironphases}
常温の純鉄\(\alpha\)-FeはBCCで強磁性、Coも常温の強磁性金属である\cite{princeton_iron,openstax_magnetism}。加熱時には約770°Cで強磁性を失い、BCCのまま常磁性となる。Curie温度1043 Kを摂氏へ換算すると約770°Cである\cite{iron_curie1954}。

その後、約912°Cで\(\gamma\)-FeのFCCへ変態する。加熱時の\(\alpha\to\gamma\)は収縮し、冷却時の逆変態では膨張する\cite{iron_transform2023}。磁気転移と結晶の同素変態は別であり、強磁性が結晶変態温度まで続くとしない。一般の熱膨張と変態による急な変化も別に考える。原子数と実際の格子定数が体積を決めるので、充填率だけで全材料の変態体積を決めない。

\subsection{被膜質量：寸法と百分率の分母}
\label{theory:material:coating}
面積\(S\)、被膜厚\(t_c\)、基材厚\(t_s\)に対し、片面被膜の質量は\(m_c=St_c\rho_c\)、基材は\(m_s=St_s\rho_s\)である。被膜の質量百分率は\(100m_c/(m_c+m_s)\)となる。厚さを同じ長さ単位にそろえ、片面・両面を図示条件に合わせる。

令和4年度基礎4-3の原典値を独自計算すると、
\[
 m_{\rm Ni}=0.50\times1.0\times10\times10^{-6}\times8900
 =0.0445\ {\rm kg},\qquad
 m_s=0.50\times1.0\times0.60\times10^{-3}\times7900
 =2.37\ {\rm kg}.
\]
したがって\(100\times0.0445/(2.37+0.0445)=1.843\ldots\%\)である。候補に対応する丸め値は\(4.5\times10^{-2}\)kg、1.8\%となる。百分率は鋼だけを分母にしない。

Niのレアメタル分類は日本の政策文脈であり、JOGMECは国際的にはベースメタルと呼ばれることも明記する\cite{jogmec_nickel2014}。分類語を普遍的な自然科学上の属性と思わず、計算条件とは別に確認する。

\subsection{材料選択：軽さ・剛性・加工性・使用環境}
\label{theory:material:selection}
軽量化でも、破断・永久変形を避ける選択と、弾性変形を抑える選択は異なる。比強度は\(\sigma_{\rm strength}/\rho\)、比弾性率は\(E/\rho\)である。対象の強度が引張・降伏のどちらかも示す。令和6年度基礎4-4はこの二つの取り違えを問う。

室温代表値では電気抵抗率はCuよりAu、AuよりFeが大きい\cite{openstax_resistivity}。純度と温度をそろえる。ステンレス鋼の一般的な定義はCr最低10.5質量\%、C最大1.2質量\%である\cite{stainless_definition2022}。したがってCr10\%未満は該当しないが、「10\%以上なら全て該当」という逆命題は導けない。

ダイヤモンドの高硬度を切削工具へ利用する例がある\cite{nims_diamond2007}。用途に応じて熱条件・被削材との化学反応も評価する。水素脆化は水素の侵入に伴う延性・破壊抵抗等の機械特性の劣化である。「物性が変化する」というだけでは範囲が広すぎ、改善や無関係な変化まで脆化とは呼ばない\cite{nims_hydrogen2024}。

硫黄快削鋼ではMnS等の介在物を利用して被削性を改善する。ただし工具寿命と仕上面粗さは同じ方向に良くなるとは限らず、脱酸・介在物の形・組織・切削速度も影響する\cite{sulfur_machinability2004}。S量だけから全鋼の被削性を保証せず、延性・溶接性など他の性能との交換条件を考える。
